\documentclass[11pt]{article}
\usepackage{fontspec}
\usepackage{xeCJK}
\usepackage{xurl}
\usepackage[a4paper,margin=25mm]{geometry}
\usepackage{amsmath,amssymb,amsthm,mathtools}
\usepackage{xcolor}
\usepackage{booktabs,array,enumitem,fancyhdr,hyperref}
\hypersetup{colorlinks=true,linkcolor=blue!50!black,urlcolor=blue!50!black}
\setlength{\emergencystretch}{2em}
\setlength{\parskip}{0.3em}
\linespread{1.16}
\allowdisplaybreaks
\newtheorem{theorem}{定理}[section]
\newtheorem{proposition}[theorem]{命题}
\newtheorem{lemma}[theorem]{引理}
\theoremstyle{definition}
\newtheorem{definition}[theorem]{定义}
\theoremstyle{remark}
\newtheorem{remark}[theorem]{说明}
\newcommand{\dd}{\,\mathrm d}
\newcommand{\ii}{\mathrm i}
\newcommand{\Rea}{\operatorname{Re}}
\newcommand{\Ima}{\operatorname{Im}}
\newcommand{\Li}{\operatorname{Li}}
\newcommand{\supp}{\operatorname{supp}}
\pagestyle{fancy}
\fancyhf{}
\fancyhead[L]{\small 学习笔记}
\fancyhead[R]{\small 初稿}
\fancyfoot[C]{\thepage}
\setlength{\headheight}{15pt}
\author{}
\date{2026年10月22日初稿}
\makeatletter
\renewcommand{\maketitle}{%
  \thispagestyle{plain}%
  \begin{center}
    {\LARGE\bfseries \@title\par}
    \vspace{0.6em}
    {\small \@date\par}
  \end{center}
  \vspace{0.6em}}
\makeatother
\newcommand{\draftnotice}{\noindent{\small
本文为供逐篇审阅的学习笔记初稿，不是新增研究成果。
文中区分直接证明与引用结果；正文不依赖本文。
本轮已作公式、条件及排版初步核查，尚待作者逐句审定。}\par\medskip}

\XeTeXlinebreaklocale "zh"
\renewcommand{\proofname}{证明}
\newcommand{\li}{\operatorname{li}}
\date{2026年10月24日初稿}
\title{黎曼猜想在素数分布中的确切位置}
\begin{document}
\maketitle

\noindent{\small 本篇说明 RH 作为数学命题所断言的内容、它能推出的素数分布结论，
以及已知证明采用更广假设或尚无由 RH 解决的问题。所引定理均注明出处；
标为“猜想”的命题均未证明。}\par\medskip

\section{问题}
在普及读物中，RH 常被描述为“质数的全部规律”，并与加密、人工智能乃至生命科学相联系。
本篇的结论是：RH 是关于 $\zeta$ 零点位置的精确命题，等价于具有最优幂指数的素数计数误差界；这里不声称对数因子已经最优。
它在素数分布理论中处于中心位置，但\emph{远不能解决素数的一切问题}。
许多基本问题需要更强的假设，另有许多问题即使 RH 与广义 RH 都成立，也没有已知的解决途径。

\paragraph{关于“蕴涵”的说明}
对尚未解决的命题 $A,B$，“$A$ 不蕴涵 $B$”无法在真值意义上断言：若 $B$ 为真，则任何 $A$ 都蕴涵 $B$。
本文所说“RH 不能解决某问题”，确切含义是：目前没有从 RH 出发的证明，并且由 RH 推出的已知结论在量上不足以得到它（参见学习笔记第 14 篇第 2 节）。

\section{素数计数的三个层次}
记 $\pi(x)$ 为不超过 $x$ 的素数个数。关于 $\pi(x)$ 的问题可以分为三个层次。
\begin{enumerate}
 \item \textbf{精确表示。} Riemann 显式公式（本书第 14 章）把 $\pi(x)$（经 $J(x)$ 与 Möbius 反演）精确地表示为对 $\zeta$ 全部非平凡零点的求和。
 这是一个对偶关系：素数完全决定零点，零点也完全决定素数。
 它把素数计数问题完整地转化为零点问题，但没有确定零点的位置；
 仅写出这一恒等式，尚未完成零点贡献与截断余项的定量控制；无条件的零点区域与计数估计仍可用来导出误差界。
 \item \textbf{主项。} 素数定理 $\pi(x)\sim\li(x)$（1896）。它依赖于 $\zeta$ 在 $\Rea s=1$ 上无零点；
 现有的误差估计 $\pi(x)-\li(x)\ll x\exp\bigl(-c(\log x)^{3/5}(\log\log x)^{-1/5}\bigr)$ 来自 Vinogradov--Korobov 型无零区域。
 \item \textbf{偏差。} $\pi(x)$ 与 $\li(x)$ 之差能有多大。这是 RH 所回答的问题：
 \begin{theorem}[von Koch 1901，引用]
 RH 成立当且仅当 $\pi(x)=\li(x)+O(\sqrt x\log x)$。
 \end{theorem}
 在 RH 下还有显式形式：$x\ge2657$ 时 $|\pi(x)-\li(x)|<\frac1{8\pi}\sqrt x\log x$（Schoenfeld 1976）。
 由 Littlewood（1914），$\pi(x)-\li(x)=\Omega_\pm\bigl(\sqrt x\log\log\log x/\log x\bigr)$，
 故 $\sqrt x$ 的量级除对数因子外不能再改进。
\end{enumerate}
因此，就 $\pi(x)$ 而言，RH 给出的误差界具有最优的幂指数 $1/2$：它约束素数在\emph{大尺度累计意义下}偏离主项的程度，而不确定误差的精确量级。
素数分布的其他问题，涉及局部间距、算术结构或加法组合，不在这一陈述之内。

\section{RH 能推出的素数分布结论}
\begin{itemize}
 \item \textbf{误差阶。} 如上节，以及 $\psi(x)=x+O(\sqrt x\log^2x)$。
 \item \textbf{相邻素数的间距。} Cramér（1920）在 RH 下证明
 $p_{n+1}-p_n=O(\sqrt{p_n}\log p_n)$；显式常数的改进见 Dudek（2015）、Carneiro--Milinovich--Soundararajan（2019）。
 无条件的最好结果是 $p_{n+1}-p_n\ll p_n^{0.525}$（Baker--Harman--Pintz 2001）。
 \item \textbf{几乎所有短区间。} Selberg（1943）在 RH 下证明：若 $f(x)\to\infty$，
 则除去一个测度意义下密度为零的集合外，对所有 $x$，区间 $(x,x+f(x)\log^2x]$ 中有素数。
\end{itemize}

\subsection{一个界限：Legendre 猜想}
Legendre 猜想断言：对每个正整数 $n$，$n^2$ 与 $(n+1)^2$ 之间有素数。
\begin{proposition}\label{prop:legendre}
由 RH 推出的间距界 $p_{n+1}-p_n=O(\sqrt{p_n}\log p_n)$ 不足以推出 Legendre 猜想。
\end{proposition}
\begin{proof}
令 $x=n^2$。区间 $(n^2,(n+1)^2)$ 的长度为 $2n+1=2\sqrt x+1$。
间距界给出的保证长度为 $C\sqrt x\log x$，当 $x$ 足够大时远大于 $2\sqrt x+1$。
因此仅凭该间距界，不能排除长为 $2\sqrt x+1$ 的区间中没有素数。
\end{proof}
这说明：即使是“相邻平方数之间有素数”这样陈述简单的问题，也不能仅由上述已知间距界得到。
更强的 Cramér 猜想 $p_{n+1}-p_n=O(\log^2p_n)$，目前同样没有仅由 RH 推出的证明。

\section{需要广义 RH 的问题}
广义 Riemann 猜想（GRH）断言全部 Dirichlet $L$ 函数（以及更一般的 $L$ 函数）的非平凡零点都在临界线上。
GRH 是比 RH 强得多的假设。普及讨论中常把 GRH 的推论归于 RH。
\begin{itemize}
 \item \textbf{等差数列中的素数。} 在 GRH 下，对 $q\le x$ 一致地有
 $\psi(x;q,a)=\frac{x}{\varphi(q)}+O(\sqrt x\log^2x)$（Davenport, \emph{Multiplicative Number Theory}, Ch.~20）。
 无条件地，Siegel--Walfisz 定理只对 $q\le(\log x)^A$ 成立；Bombieri--Vinogradov 定理在对 $q$ 求平均的意义下给出与 GRH 相当的结论。
 RH 本身对此没有直接作用。
 \item \textbf{最小模素数。} 模 $q$ 剩余类 $a$ 中的最小素数：无条件地 $\ll q^L$（Linnik；目前 $L\le5$，Xylouris 2011）；
 在 GRH 下 $\ll q^{2+\varepsilon}$。
 \item \textbf{Artin 原根猜想。} 非平方数 $a\ne-1$ 是无穷多个素数的原根。
 Hooley（1967）在 Kummer 域 Dedekind zeta 函数的 GRH 下证明了它；
 无条件地，Heath-Brown（1986）证明至多有两个素数例外。
 \item \textbf{三素数定理。} 每个充分大的奇数是三个素数之和。Hardy 与 Littlewood（1923）在 GRH 下证明；
 Vinogradov（1937）无条件地证明了充分大的情形；Helfgott（2013）完成了全部奇数 $\ge7$ 的证明。
 这是一个“在 GRH 下先得证、后来无条件解决”的例子。
\end{itemize}

\section{尚无由 RH 或 GRH 解决的若干问题}
以下问题，目前没有从 RH 或 GRH 出发的解决途径；其中的进展大多来自筛法与组合方法。
\begin{itemize}
 \item \textbf{孪生素数猜想。} 无穷多对素数相差 $2$。
 Zhang（2013）证明了有界间距，Maynard 与 Tao（2015）及 Polymath 项目把界缩小到 $246$；
 这些结果使用 Bombieri--Vinogradov 定理与筛法，并不依赖 RH。
 若假设 Elliott--Halberstam 猜想（一个与 GRH 不同的分布猜想），间距可缩小到 $12$（Maynard）。
 \item \textbf{偶数 Goldbach 猜想。} 每个大于 $2$ 的偶数是两个素数之和。
 相关的 GRH 条件性结果给出例外集估计：不超过 $x$ 的例外偶数至多 $O(x^{1/2+\varepsilon})$ 个（Hardy--Littlewood 1924）；
 这一估计没有排除全部例外，目前也没有仅由 GRH 推出猜想本身的证明。
 \item \textbf{素数间距的极端值。} 大间距的下界（Ford--Green--Konyagin--Maynard--Tao 2018）是无条件的组合结果；
 Cramér 猜想尚不能由上述 RH 条件下的已知间距界推出（第 3.1 节）。
 \item \textbf{多项式取素数值。} 例如 $n^2+1$ 是否无穷多次为素数（Landau 问题之一）；目前没有仅由 RH 解决这一问题的证明。
 \item \textbf{素数中的结构。} Green--Tao 定理（2008）：素数中有任意长的等差数列。它的证明不使用 RH。
 \item \textbf{素数竞赛。} 模 $4$ 余 $3$ 的素数“通常”多于余 $1$ 的素数（Chebyshev 偏差）。
 Rubinstein 与 Sarnak（1994）对偏差的精确刻画需要 GRH \emph{加上}零点虚部线性无关的假设（LI）；GRH 本身不够。
\end{itemize}

\section{汇总}
\begin{center}
\small
\renewcommand{\arraystretch}{1.25}
\begin{tabular}{>{\raggedright\arraybackslash}p{0.38\textwidth}>{\raggedright\arraybackslash}p{0.52\textwidth}}
\toprule
问题 & 与 RH 的关系\\
\midrule
$\pi(x)$ 的精确表示 & 已由显式公式给出，与 RH 无关\\
素数定理（主项） & 已证明，依赖 $\Rea s=1$ 上无零点\\
$\pi(x)-\li(x)=O(\sqrt x\log x)$ & \textbf{等价于 RH}，幂指数最优\\
相邻素数间距 $O(\sqrt p\log p)$ & 由 RH 推出\\
Legendre 猜想、Cramér 猜想 & 尚不能由上述已知间距界推出\\
等差数列中素数的误差、最小模素数 & 上述条件性证明采用 GRH\\
Artin 原根猜想 & Hooley 的条件性证明采用相应 GRH\\
三素数定理 & 曾需 GRH，现已无条件证明\\
孪生素数、偶数 Goldbach & RH 与 GRH 均无已知途径\\
素数竞赛的精确偏差 & 需要 GRH 与 LI\\
素数中的长等差数列 & 已证明，不依赖 RH\\
\bottomrule
\end{tabular}
\end{center}

\section{与计算和应用的关系}
\begin{itemize}
 \item \textbf{加密。} RSA 等方案的安全性基于大整数分解的计算困难。RH 是关于素数分布的命题，
 不提供分解算法；它的证明或否证都不会直接改变这些方案的安全性。
 \item \textbf{素性检验。} Miller（1976）的确定性素性检验以 GRH 为前提（Bach 1990 给出显式界）。
 AKS 算法（2004）无条件地给出多项式时间的确定性检验，使这一依赖不再必要。
 \item \textbf{计算 $\pi(x)$。} 实际计算 $\pi(x)$ 使用组合方法（Meissel--Lehmer、Lagarias--Miller--Odlyzko）
 或基于 $\zeta$ 积分的解析方法，都不以 RH 为前提。
 \item \textbf{物理。} 零点间距的统计与随机矩阵理论（GUE）一致，后者也描述重原子核能级的统计（学习笔记第 15 篇）。
 这是统计规律上的深刻对应，并不意味着 RH 与生物或医学问题有关。
\end{itemize}

\section{有限检验与猜想的证据}
RH 等价于一个 $\Pi_1$ 型算术命题（学习笔记第 13 篇）：若它不成立，存在有限的反例，可以由计算严格证实。
在这一意义上它是可证伪的。高度 $3\cdot10^{12}$ 以下的全部零点已被严格证明位于临界线上（第 12 篇），
零点比例、零点密度、函数域类比与随机矩阵统计也都与 RH 相容（第 15 篇）。
这些结果从不同方面支持对 RH 的研究，但各自的适用范围不同，不能合成为覆盖全部零点的证明。

但数学的标准不同。实验科学相信一条定律，前提是检验覆盖了它起作用的范围；
而控制零点局部行为的量随 $\log\log T$ 变化，现有数据只覆盖 $\log\log T\le3.4$（第 12 篇），
数论中也有 Pólya、Mertens 猜想与 $\pi(x)<\li(x)$ 在大范围验证后失效的先例。
因此，大范围检验及其他理论证据支持 RH，却没有改变它仍是未证猜想的地位。

\section{为何 RH 仍居中心地位}\label{sec:center}
前几节说明了 RH 不能做什么。它的中心地位来自以下几点，而不是来自“解决素数的一切问题”。

\subsection{误差阶的含义：平方根相消}
由 Littlewood 定理（第 2 节），$\pi(x)-\li(x)$ 的量级不能小于 $\sqrt x$（除对数因子外）；
RH 等价于上界 $\pi(x)-\li(x)=O(\sqrt x\log x)$。上下界的对数因子并不相同；这里所谓平方根尺度只指幂指数 $1/2$。
这一量级恰与独立随机符号部分和的量级相同：独立 $\pm1$ 序列的部分和几乎必然为 $O(N^{1/2+\varepsilon})$
（学习笔记第 15 篇命题 4.1），而 RH 等价于 $M(x)=\sum_{n\le x}\mu(n)=O(x^{1/2+\varepsilon})$（第 14 篇第 3.2 节）。
在这个意义上，RH 说的是：素数与 Möbius 函数在大尺度上的累计偏差，与随机序列一样小。

“平方根相消”是解析数论中反复出现的原则，指数和、特征和、Kloosterman 和都以它为理想情形。
RH 是这一原则在素数上的基本形式。

\subsection{已证明的类似命题是实际使用的工具}
有限域上的类似命题已被证明（Weil 1948，Deligne 1974；学习笔记第 15 篇第 3.1 节），并给出大量平方根相消的估计。
例如，对素数 $p$ 与 $p\nmid ab$，Kloosterman 和满足
\[
 \Bigl|\sum_{x=1}^{p-1}e^{2\pi\ii(ax+b\bar x)/p}\Bigr|\le2\sqrt p
\]
（Weil 1948，$\bar x$ 为 $x$ 模 $p$ 的逆）。这类估计是模形式、筛法与素数分布研究中的基本输入。
第 5 节所说孪生素数问题上的进展（Zhang 2013）不依赖经典的 RH，却用到了 Deligne 证明的有限域上的 RH 所给出的指数和估计。
可见 RH 型命题的作用远超素数计数：它们在已被证明的场合已经成为工作工具。

\subsection{条件性定理}
第 3、4 节所列的许多结论在 RH 或 GRH 下已经证明。数论文献中以此为前提的结果数以百计。
RH 及其推广若被证明，以相应假设为前提的结论便可无条件使用；若假设被否证，已证明的条件命题仍成立，但不能再凭该假设得到无条件结论。
RH 的分量在很大程度上来自它所支撑的这一整体。

\subsection{困难本身的意义}
在有限域的情形，$\zeta$ 函数的零点是 Frobenius 映射在上同调上作用的特征值的倒数，
RH 归结为这些特征值的模，证明依赖这一几何结构。对 $\mathbb Q$ 上的 $\zeta$，目前不知道对应的算子或几何对象；
Hilbert--Pólya 设想（零点是某个自伴算子的谱）只是猜测（学习笔记第 15 篇第 3.2 节）。
寻找可能解释零点位置的谱或几何结构，是有价值的探索方向；如果发现这样的结构，其作用可能超出 RH 本身。
这是一种研究设想，并不意味着已经知道 RH 必须依靠新的结构才能证明。

\section{结论}
\begin{enumerate}
 \item RH 是一个精确的数学命题：$\zeta$ 的非平凡零点都在临界线上；等价地，$\pi(x)-\li(x)=O(\sqrt x\log x)$。
 \item 显式公式已经把素数计数完全转化为零点问题；RH 回答的是这一转化之后的误差阶。
 \item RH 不能解决素数的一切问题：局部间距（Legendre、Cramér 猜想）、等差数列中的分布（需要 GRH）、
 加法与模式问题（孪生素数、Goldbach）都不在它的范围内。
 \item RH 的中心地位来自四个方面：它表达了素数所能达到的最大均匀性（平方根相消）；
 它已被证明的有限域类似命题是实际使用的工具；大量条件性定理以它为前提；
 对潜在谱与几何结构的探索可能带来新的理解（第~\ref{sec:center}~节）。这是方法与结构上的重要性，而不是对素数分布的全面解答。
\end{enumerate}

\section*{参考}
\begin{itemize}
 \item H.\ von Koch, Sur la distribution des nombres premiers, \emph{Acta Math.} 24 (1901)；
 L.\ Schoenfeld, Sharper bounds for the Chebyshev functions $\theta(x)$ and $\psi(x)$ II, \emph{Math.\ Comp.} 30 (1976)。
 \item H.\ Cramér, Some theorems concerning prime numbers, \emph{Ark.\ Mat.\ Astr.\ Fys.} 15 (1920)（RH 下的间距界）；
 H.\ Cramér, On the order of magnitude of the difference between consecutive prime numbers, \emph{Acta Arith.} 2 (1936)（Cramér 猜想）；
 A.\ Selberg, On the normal density of primes in small intervals, and the difference between consecutive primes, \emph{Arch.\ Math.\ Naturvid.} 47 (1943)。
 \item R.\ C.\ Baker, G.\ Harman, J.\ Pintz, The difference between consecutive primes II, \emph{Proc.\ London Math.\ Soc.} 83 (2001)。
 \item H.\ Davenport, \emph{Multiplicative Number Theory}, 3rd ed., Springer, 2000（等差数列中的素数、Bombieri--Vinogradov 定理）。
 \item C.\ Hooley, On Artin's conjecture, \emph{J.\ reine angew.\ Math.} 225 (1967)；
 D.\ R.\ Heath-Brown, Artin's conjecture for primitive roots, \emph{Quart.\ J.\ Math.} 37 (1986)。
 \item H.\ A.\ Helfgott, The ternary Goldbach problem, 2013/2015, \url{https://arxiv.org/abs/1501.05438}。
 \item Y.\ Zhang, Bounded gaps between primes, \emph{Ann.\ of Math.} 179 (2014)；
 J.\ Maynard, Small gaps between primes, \emph{Ann.\ of Math.} 181 (2015)。
 \item B.\ Green, T.\ Tao, The primes contain arbitrarily long arithmetic progressions, \emph{Ann.\ of Math.} 167 (2008)。
 \item M.\ Rubinstein, P.\ Sarnak, Chebyshev's bias, \emph{Experiment.\ Math.} 3 (1994)。
 \item M.\ Agrawal, N.\ Kayal, N.\ Saxena, PRIMES is in P, \emph{Ann.\ of Math.} 160 (2004)。
 \item A.\ Weil, On some exponential sums, \emph{PNAS} 34 (1948)；P.\ Deligne, La conjecture de Weil I, \emph{Publ.\ Math.\ IHÉS} 43 (1974)。
 \item 本书第 14、15 章；学习笔记第 12--15 篇。
\end{itemize}
\end{document}
